%%% FIELD proof-learned-projection-convergence
Fix a fold \(k\) and suppress its construction and calibration superscripts.  Define
\[
 \Delta h_{aj}=\widehat h_{aj}^{(-k,A)}-h^0_{aj},\qquad
 \Delta q_{a\ell}=\widehat q_{a\ell}^{(-k,A)}-q^0_{a\ell}.
\]
Direct expansion of the signal in \(\mathrm{def:score-family}\) gives the exact cell identity
\[
 S_a(\widehat h_{aj},\widehat q_{a\ell})-S_a(h^0_{aj},q^0_{a\ell})
 =(1-I_aq^0_{a\ell})\Delta h_{aj}
 +I_a(Y-h^0_{aj})\Delta q_{a\ell}
 -I_a\Delta q_{a\ell}\Delta h_{aj}.                                      \tag{1}
\]
All products in (1) are square-integrable by \(\mathrm{ass:heldout-signal-integrability}\) and the fourth-moment hypotheses below.  By H\"older's inequality, \(\mathrm{ass:outcome-learning-rates}\), \(\mathrm{ass:treatment-learning-rates}\), and \(\mathrm{ass:moment-envelope-growth}\), uniformly over retained cells and folds,
\[
\begin{aligned}
 \|S_a(\widehat h_{aj},\widehat q_{a\ell})-S_a(h^0_{aj},q^0_{a\ell})\|_2
 &\le M_{q,n}\|\Delta h_{aj}\|_4+M_{y,n}\|\Delta q_{a\ell}\|_4
       +\|\Delta h_{aj}\|_4\|\Delta q_{a\ell}\|_4\\
 &=O_p(d^S_{a,j\ell,n}).                                                    \tag{2}
\end{aligned}
\]
Each pure or four-corner coordinate is a signed sum of at most four cell signals.  Thus the triangle inequality, followed by the definition in \(\mathrm{def:inference-gauges}\), yields
\[
 \|\widehat D-D^0\|_{L^2(P;\ell^2)}^2
 =\sum_{c=1}^{J_n}\|\widehat D_c-D_c^0\|_2^2
 =O_p\!\left(\sum_{c=1}^{J_n}d_{c,n}^2\right).
\]
Because \(K_{\mathrm{fold}}\) is fixed, the same bound holds for the maximum over folds.  Consequently,
\[
 \max_k\|\widehat D_{J_n}^{(-k,A)}-D^0_{J_n}\|_{L^2(P;\ell^2)}
 =O_p(\delta_{D,n}),\qquad
 \delta_{D,n}\le \sqrt{J_n}\max_{c\le J_n}d_{c,n}.                       \tag{3}
\]
Applying (2) to the two base cells gives, for
\[
 B^0=\sum_{a=0}^1t_aS_a(h^0_{a0},q^0_{a0})=\theta(P)+\phi_0,
\]
the further bound
\[
 \max_k\|\widehat B^{(-k,A)}-B^0\|_2=O_p(d_{0,n}).                         \tag{4}
\]

Put
\[
 G=E[D^0D^{0\top}],\quad c=E[D^0B^0],\quad
 \widetilde G=E[\widehat D\widehat D^\top\mid\mathcal D_{A,k}],\quad
 \widetilde c=E[\widehat D\widehat B\mid\mathcal D_{A,k}].
\]
The population envelope in \(\mathrm{ass:moment-envelope-growth}\) implies
\(\|D^0\|_{L^2(P;\ell^2)}\le\sqrt{J_n}B_{D,n}\) and \(\|B^0\|_2\le B_{0,n}\).  The Hilbert-space Cauchy--Schwarz inequality, (3), and (4) therefore give
\[
 \|\widetilde G-G\|_{\mathrm{op}}
 \le \|\widehat D-D^0\|_{L^2(P;\ell^2)}
       \{\|\widehat D\|_{L^2(P;\ell^2)}+\|D^0\|_{L^2(P;\ell^2)}\}
 =O_p\{2\sqrt{J_n}B_{D,n}\delta_{D,n}+\delta_{D,n}^2\},                  \tag{5}
\]
and
\[
 \|\widetilde c-c\|_2
 =O_p\!\left\{B_{0,n}\delta_{D,n}
 +(\sqrt{J_n}B_{D,n}+\delta_{D,n})d_{0,n}\right\}.                       \tag{6}
\]
The established \(\mathrm{lem:conditional-calibration-concentration}\) supplies calibration deviations of orders \(\eta_{G,n}\) and \(\eta_{c,n}\).  Combining that lemma with (5)--(6) and \(\mathrm{def:inference-gauges}\) proves
\[
 \max_k\|\widehat G_k^C-G\|_{\mathrm{op}}=O_p(\zeta_{G,n}),\qquad
 \max_k\|\widehat c_k^C-c\|_2=O_p(\zeta_{c,n}).                          \tag{7}
\]
In particular, the sharpened conditional rates used here are
\[
 \eta_{G,n}=B_{D,n}^2\left\{\sqrt{J_n/n_{C,\min}}+J_n/n_{C,\min}\right\},
 \qquad
 \eta_{c,n}=B_{D,n}B_{0,n}\sqrt{J_n/n_{C,\min}}.
\]
Thus, when the displayed envelopes are bounded, the calibration part is consistent under \(J_n/n_{C,\min}\to0\); no logarithmic dimension factor is required.

Let \(\Pi_R=\Pi_{R,J_n}\) project onto \(\operatorname{range}(G)\), and write
\(\delta_k=\|\widehat G_k^C-G\|_{\mathrm{op}}\).  By \(\mathrm{ass:gram-conditioning}\), every positive eigenvalue of \(G\) is at least \(\kappa_{J_n}\), \(0<\tau_n<\kappa_{J_n}/2\), and \(\delta_k=o_p(\tau_n)\) uniformly over the fixed number of folds.  The variational inequality
\[
 |u^\top(\widehat G_k^C-G)u|\le\delta_k\|u\|_2^2                         \tag{8}
\]
shows that every empirical eigenvalue paired with the population zero spectrum is at most \(\delta_k<\tau_n\), whereas the remaining eigenvalues are at least \(\kappa_{J_n}-\delta_k>\tau_n\), with probability tending to one.  Hence the thresholded operator
\[
 \widehat G_{k,\tau}=\widehat\Pi_k^C\widehat G_k^C\widehat\Pi_k^C
\]
has the same rank as \(G\), reduced minimum modulus at least \(\kappa_{J_n}/2\), and
\(\|\widehat G_{k,\tau}-G\|_{\mathrm{op}}\le2\delta_k\).  Applying the established \(\mathrm{lem:operator-kernel-perturbation}\) to these finite-dimensional self-adjoint operators and taking complements of their kernel projections yields
\[
 \max_k\|\widehat\Pi_k^C-\Pi_R\|_{\mathrm{op}}
 =O_p(\zeta_{G,n}/\kappa_{J_n})=O_p(\eta_{\Pi,n}).                         \tag{9}
\]

The population least-squares coefficient satisfies
\[
 Gb_{J_n}=c,\qquad b_{J_n}\in\operatorname{range}(G).                     \tag{10}
\]
Indeed, if \(v\in\ker G\), then \(E[(v^\top D^0)^2]=v^\top Gv=0\), so \(v^\top D^0=0\) almost surely and \(v^\top c=E[v^\top D^0B^0]=0\).  Hence \(c\in\operatorname{range}(G)\), which proves (10) on the nonzero Gram eigenspace specified in \(\mathrm{def:population-affine-sieve}\).  From \(\mathrm{def:affine-ensemble-estimator}\), the cutoff-ridge coefficient obeys
\[
 (\widehat G_k^C+\rho_n I)\widehat b=\widehat\Pi_k^C\widehat c_k^C,
 \qquad \widehat b=\widehat\Pi_k^C\widehat b.                            \tag{11}
\]
Subtracting (10) from (11), using \(G(I-\Pi_R)=0\), and projecting onto \(\operatorname{range}(G)\) gives
\[
 G\Pi_R(\widehat b-b_{J_n})
 =\Pi_R\{(G-\widehat G_k^C)\widehat b+(\widehat c_k^C-c)
 -\rho_n\widehat b-(I-\widehat\Pi_k^C)\widehat c_k^C\}.                 \tag{12}
\]
Since \(c=\Pi_Rc\), (7) and (9) imply
\[
 \|(I-\widehat\Pi_k^C)\widehat c_k^C\|_2
 \le \|\widehat c_k^C-c\|_2+\|(\Pi_R-\widehat\Pi_k^C)c\|_2
 =O_p\{\zeta_{c,n}+\eta_{\Pi,n}\|c\|_2\}.                             \tag{13}
\]
For each unit vector \(u\), Cauchy--Schwarz gives
\(
 |u^\top c|\le\{E[(u^\top D^0)^2]\}^{1/2}\|B^0\|_2
\), and therefore
\[
 \|c\|_2\le\sqrt{\Lambda_{J_n}}B_{0,n}.                                \tag{14}
\]
Invert \(G\) only on its positive eigenspace.  Equations (7), (12)--(14), the lower eigenvalue bound in \(\mathrm{ass:gram-conditioning}\), and the coefficient envelope in \(\mathrm{ass:affine-weight-envelope}\) prove
\[
 \max_k\|\Pi_R(\widehat b-b_{J_n})\|_2
 =O_p\!\left[
 \frac{\zeta_{c,n}+(\zeta_{G,n}+\rho_n)L_n
 +\sqrt{\Lambda_{J_n}}B_{0,n}\eta_{\Pi,n}}
 {\kappa_{J_n}}
 \right].                                                                  \tag{15}
\]
This retains the learned coefficient multiplying the Gram perturbation and the ridge term.  Furthermore,
\[
 \|v^\top D^0\|_2=\|G^{1/2}v\|_2
 \le\sqrt{\Lambda_{J_n}}\|\Pi_Rv\|_2,
 \qquad
 \|\widehat b^\top(\widehat D-D^0)\|_2\le L_n\delta_{D,n}.
\]
Together with (3) and (15), these inequalities give
\[
\begin{aligned}
 &\max_k\|\widehat b^{(-k,C\mid A)\top}\widehat D_{J_n}^{(-k,A)}
       -b_{J_n}^\top D^0_{J_n}\|_2\\
 &\quad=O_p\!\left[L_n\delta_{D,n}
 +\frac{\sqrt{\Lambda_{J_n}}}{\kappa_{J_n}}
 \left\{\zeta_{c,n}+(\zeta_{G,n}+\rho_n)L_n
 +\sqrt{\Lambda_{J_n}}B_{0,n}\eta_{\Pi,n}\right\}\right].             \tag{16}
\end{aligned}
\]

It remains to identify the adjusted signal.  Collecting the coefficient of every bridge cell in the deterministic coefficient-to-weight map in \(\mathrm{def:affine-ensemble-estimator}\) gives, for each arm,
\[
 \sum_{j=0}^{J_n^h}\sum_{\ell=0}^{J_n^q}\widehat\lambda_{a,j\ell}=1,
 \qquad
 \widehat\psi=\widehat B-\widehat b^\top\widehat D.                     \tag{17}
\]
The identical algebra applies to the population coefficient.  By the established \(\mathrm{lem:normal-affine-identity}\), every coordinate of \(D^0\) belongs to \(\mathcal C_P\).  Orthogonal decomposition therefore gives, for every \(b\in\mathbb R^{J_n}\),
\[
 \|\phi_0-b^\top D^0\|_2^2
 =\|\phi_0-\Pi_{\mathcal C_P}\phi_0\|_2^2
  +\|\Pi_{\mathcal C_P}\phi_0-b^\top D^0\|_2^2.                         \tag{18}
\]
The normal equation (10) shows that \(b_{J_n}\) minimizes the \(b\)-dependent term in (18), whose norm is therefore \(\alpha_{J_n}\).  Under the hypotheses of the established \(\mathrm{thm:quotient-efficiency}\),
\(\phi_{\mathrm{eff}}=\phi_0-\Pi_{\mathcal C_P}\phi_0\).  Since \(B^0=\theta(P)+\phi_0\), it follows that
\[
 \|B^0-b_{J_n}^\top D^0-\theta(P)-\phi_{\mathrm{eff}}\|_2
 =\alpha_{J_n}.                                                             \tag{19}
\]
Combining (4), (16), (17), and (19) gives exactly
\[
 \max_k\|\widehat\psi^{(-k,C\mid A)}-\theta(P)-\phi_{\mathrm{eff}}\|_2
 =O_p(\Gamma_n).
\]
Finally, \(\mathrm{ass:projection-rate}\) states \(\Gamma_n=o_p(1)\); its constituent approximation, conditioning, ridge, and learning requirements are supplied by the other displayed hypotheses.  This proves every assertion of the lemma.
%%% FIELD proof-affine-ensemble-inference
We first prove the pointwise assertions.  Let \(P_{n,k}\) denote the empirical measure on \(\mathcal I_k\), and define the random function
\[
 f_k=\widehat\psi^{(-k,C\mid A)}-\theta(P)-\phi_{\mathrm{eff}}.
\]
The proved \(\mathrm{lem:learned-projection-convergence}\), with its nested-split hypotheses, gives
\[
 \max_{k\le K_{\mathrm{fold}}}\|f_k\|_{L^2(P)}=o_p(1).                   \tag{20}
\]
By \(\mathrm{ass:training-measurability}\), \(f_k\) is measurable with respect to \(\mathcal D_{A,k}\vee\mathcal D_{C,k}\), while the observations in \(\mathcal I_k\) are independent of that sigma-field under \(\mathrm{ass:iid-sampling}\).  Conditional square-integrability follows from \(\mathrm{ass:heldout-signal-integrability}\).  Since all fold sizes are proportional to \(n\) and \(K_{\mathrm{fold}}\) is fixed, conditional variance calculation yields
\[
 E\!\left[
 \left.\left\{\frac{|\mathcal I_k|}{\sqrt n}(P_{n,k}-P)f_k\right\}^2
 \right|\mathcal D_{A,k}\vee\mathcal D_{C,k}\right]
 \le C\|f_k\|_2^2=o_p(1).                                                  \tag{21}
\]
Conditional Chebyshev inequality and a finite union bound imply
\[
 \sum_{k=1}^{K_{\mathrm{fold}}}\frac{|\mathcal I_k|}{\sqrt n}
 (P_{n,k}-P)f_k=o_p(1).                                                      \tag{22}
\]

We next bound the conditional drift.  For any learned cell, (1), outcome-bridge solvability, and treatment-bridge solvability give
\[
 P\{S_a(\widehat h,\widehat q)-S_a(h^0,q^0)\}
 =-E[I_a(\widehat q-q^0)(\widehat h-h^0)].                                 \tag{23}
\]
Indeed, treating the learned functions as fixed after conditioning on their training sigma-field,
\[
 E[(\widehat h-h^0)(1-I_aq^0)]
 =E[(\widehat h-h^0)(1-T_a^*q^0)]=0
\]
because \(q^0\in\mathcal Q_a(P)\), and
\[
 E[I_a(\widehat q-q^0)(Y-h^0)]
 =E[I_a(\widehat q-q^0)(g_a-T_ah^0)]=0
\]
because \(h^0\in\mathcal H_a(P)\).  Conditioning the last expectation in (23), first on \((A,Z,X)\) and then on \((W,X)\), and applying Cauchy--Schwarz in the corresponding Hilbert spaces gives both bounds
\[
 |E[I_a\Delta q\Delta h]|
 \le\|T_a\Delta h\|_{\mathcal G_a}\|\Delta q\|_{\mathcal G_a},
 \qquad
 |E[I_a\Delta q\Delta h]|
 \le\|\Delta h\|_{\mathcal H}\|T_a^*\Delta q\|_{\mathcal H}.         \tag{24}
\]
The affine weights sum to one within each arm by (17), and their total variation is at most \(L_n\) by \(\mathrm{ass:affine-weight-envelope}\).  Hence (23)--(24), \(\mathrm{ass:outcome-residual-rate}\), \(\mathrm{ass:adjoint-residual-rate}\), and the \(L^2\) parts of \(\mathrm{ass:outcome-learning-rates}\) and \(\mathrm{ass:treatment-learning-rates}\) prove
\[
 \max_k|P\widehat\psi^{(-k,C\mid A)}-\theta(P)|
 =O_p\!\left(L_n\min\{r_{h,n}e_{q,2,n},r_{q,n}e_{h,2,n}\}\right)
 =o_p(n^{-1/2}),                                                             \tag{25}
\]
where the last equality is precisely \(\mathrm{ass:primitive-mixed-bias-rate}\).

The folds partition the sample, \(P\phi_{\mathrm{eff}}=0\) by \(\mathrm{thm:quotient-efficiency}\), and therefore (22)--(25) give
\[
\begin{aligned}
 \sqrt n\{\widehat\theta_n-\theta(P)\}
 &=\frac1{\sqrt n}\sum_{i=1}^n\phi_{\mathrm{eff}}(O_i)
 +\sum_k\frac{|\mathcal I_k|}{\sqrt n}(P_{n,k}-P)f_k
 +\sum_k\frac{|\mathcal I_k|}{\sqrt n}Pf_k\\
 &=\frac1{\sqrt n}\sum_{i=1}^n\phi_{\mathrm{eff}}(O_i)+o_p(1).          \tag{26}
\end{aligned}
\]
Also by \(\mathrm{thm:quotient-efficiency}\),
\(E[\phi_{\mathrm{eff}}^2]=V_{\mathrm{eff}}(P)\in(0,\infty)\).  For every \(\epsilon>0\), integrability and dominated convergence give
\[
 E[\phi_{\mathrm{eff}}^2\mathbf 1\{|\phi_{\mathrm{eff}}|>\epsilon\sqrt n\}]
 \longrightarrow0.
\]
Thus \(\mathrm{lem:lindeberg-feller-central-limit}\), applied to
\(X_{ni}=n^{-1/2}\phi_{\mathrm{eff}}(O_i)\), implies
\[
 n^{-1/2}\sum_{i=1}^n\phi_{\mathrm{eff}}(O_i)
 \Rightarrow N\{0,V_{\mathrm{eff}}(P)\}.                                  \tag{27}
\]
For completeness, (26) and (27) have the same limit without invoking any additional result: for every continuity point \(x\) of the normal distribution and every \(\delta>0\), the distribution function of the left side of (26) is bounded above and below by those of the leading sum at \(x+\delta\) and \(x-\delta\), plus the probability that the remainder exceeds \(\delta\); first let \(n\to\infty\), then \(\delta\downarrow0\).  This proves the asserted pointwise Gaussian limit.

For variance estimation, (20), conditional Markov inequality, and the same foldwise conditioning imply
\[
 \frac1n\sum_{i=1}^nf_{k(i)}(O_i)^2=o_p(1).                                \tag{28}
\]
The weak law for \(\phi_{\mathrm{eff}}^2\) follows under only its finite first moment: for fixed \(M\), Chebyshev's inequality applies to the bounded variables \(\phi_{\mathrm{eff}}^2\wedge M\), and Markov's inequality controls the empirical remainder by
\(E[\phi_{\mathrm{eff}}^2\mathbf 1\{\phi_{\mathrm{eff}}^2>M\}]\), which vanishes as \(M\to\infty\).  Hence
\[
 \frac1n\sum_{i=1}^n\phi_{\mathrm{eff}}(O_i)^2
 \longrightarrow_pV_{\mathrm{eff}}(P).                                   \tag{29}
\]
Empirical Cauchy--Schwarz, (28), and (29) make the cross term negligible.  Moreover, (26)--(27) imply \(\widehat\theta_n-\theta(P)=O_p(n^{-1/2})\).  Using the definition of \(\widehat V_n\), we obtain
\[
\begin{aligned}
 \widehat V_n
 &=\frac1n\sum_{i=1}^n\{\widehat\psi_i-\theta(P)\}^2
   -\{\widehat\theta_n-\theta(P)\}^2\\
 &\longrightarrow_pV_{\mathrm{eff}}(P).                                   \tag{30}
\end{aligned}
\]
Since the limit is positive, division in (30) proves
\(\widehat V_n/V_{\mathrm{eff}}(P)\to_p1\).  Combining (26)--(27) with this ratio convergence, using the same distribution-function sandwich as above, shows that the studentized statistic converges to \(N(0,1)\).  Continuity of its distribution at \(\pm z_{1-\alpha/2}\) proves
\[
 P\{\theta(P)\in\mathrm{CI}_{1-\alpha}\}\longrightarrow1-\alpha.
\]

We now prove the locally uniform extension, without resolving \(\mathrm{oeq:score-lifting-completeness}\).  Let
\[
 \Xi_n=\{(s,t):s\in\mathcal S,\ |t|\le T\},\qquad
 Q_{n,\xi}=P_{n,t,s}\quad\text{for }\xi=(s,t),
\]
and abbreviate
\[
 \theta_{n,\xi}=\theta(Q_{n,\xi}),\qquad
 \phi_{n,\xi}=\phi_{\mathrm{eff}}(Q_{n,\xi}),\qquad
 V_{n,\xi}=E_{Q_{n,\xi}}[\phi_{n,\xi}^2].
\]
The theorem's uniform hypotheses expressly impose the hypotheses of \(\mathrm{thm:quotient-efficiency}\), including tangent equality and pathwise product validity, at every \(Q_{n,\xi}\).  Thus that theorem identifies \(\phi_{n,\xi}\) as the canonical gradient at each such law.  This is a premise at the local laws, not a proof of the unresolved score-lifting equality.

Choose an arbitrary deterministic sequence \(\xi_n\in\Xi_n\) and put \(Q_n=Q_{n,\xi_n}\).  Equations (1)--(19) apply under \(Q_n\): every bridge-learning, multiplier, sieve-approximation, calibration-concentration, positive-eigenvalue, cutoff, coefficient, weight, and projection-gauge premise used there is assumed uniformly over \(\Xi_n\).  Therefore
\[
 \max_k\|\widehat\psi_{n,k}-\theta(Q_n)-\phi_{n,\xi_n}\|_{L^2(Q_n)}
 =o_{Q_n}(1).                                                               \tag{31}
\]
The uniformly root-\(n\)-small mixed-bias gauge likewise makes the analogue of (25) \(o_{Q_n}(n^{-1/2})\).  Conditional evaluation-fold Chebyshev under \(Q_n\), exactly as in (21)--(26), then proves
\[
 \sqrt n\{\widehat\theta_n-\theta(Q_n)\}
 =\frac1{\sqrt n}\sum_{i=1}^n\phi_{n,\xi_n}(O_i)+o_{Q_n}(1).              \tag{32}
\]
Because \(\xi_n\) was arbitrary, (32) holds uniformly over \(\Xi_n\).  Indeed, failure of uniformity would provide \(\epsilon,\delta>0\), a subsequence, and choices \(\xi_n\) on that subsequence for which the remainder probability is at least \(\delta\), contradicting the sequencewise argument.

To compare the varying gradients in one Hilbert space, define the canonical likelihood-ratio isometry
\[
 \mathsf J_{n,\xi}f=f\sqrt{dQ_{n,\xi}/dP}:L^2(Q_{n,\xi})\longrightarrow L^2(P).
\]
It is an isometry because
\(\|\mathsf J_{n,\xi}f\|_{L^2(P)}^2=E_{Q_{n,\xi}}[f^2]\).  The assumed uniform transported \(L^2(P)\) convergence gives
\[
 \sup_{\xi\in\Xi_n}|V_{n,\xi}-V_{\mathrm{eff}}(P)|\longrightarrow0.       \tag{33}
\]
Every \(Q_{n,\xi}\) lies in \(\mathcal V_P\), so its likelihood ratio is between \(1-\varepsilon\) and \(1+\varepsilon\).  Hence, for every fixed \(\epsilon>0\),
\[
\begin{aligned}
 &E_{Q_{n,\xi}}[\phi_{n,\xi}^2
       \mathbf 1\{|\phi_{n,\xi}|>\epsilon\sqrt n\}]\\
 &\quad\le E_P[(\mathsf J_{n,\xi}\phi_{n,\xi})^2
       \mathbf 1\{|\mathsf J_{n,\xi}\phi_{n,\xi}|
       >\epsilon\sqrt{1-\varepsilon}\sqrt n\}],
                                                                                \tag{34}
\end{aligned}
\]
which tends to zero uniformly by uniform integrability of the transported squared family; this also verifies the theorem's displayed uniform Lindeberg requirement.  For every sequence \(\xi_n\), \(\mathrm{lem:lindeberg-feller-central-limit}\), (33), and (34) give, under \(Q_{n,\xi_n}^n\),
\[
 n^{-1/2}\sum_{i=1}^n\phi_{n,\xi_n}(O_i)
 \Rightarrow N\{0,V_{\mathrm{eff}}(P)\}.                                  \tag{35}
\]
The preceding contradiction argument upgrades (35) to uniform convergence in bounded-Lipschitz distance over \(\Xi_n\).  Combining (32) and (35) proves the asserted uniform centered Gaussian limit.

The variance argument is also uniform.  Along every sequence \(\xi_n\), conditional Markov inequality applied to (31) gives the analogue of (28).  For the leading squares, truncate \(\phi_{n,\xi_n}^2\) at \(M\).  Chebyshev's inequality makes the centered average of the truncated variables vanish because its variance is at most
\(M E_{Q_n}[\phi_{n,\xi_n}^2]/n\), which is uniformly \(O(M/n)\) by (33).  Uniform integrability in (34) controls both the empirical and population tails as \(M\to\infty\).  The arbitrary-sequence argument thus yields, for every \(\epsilon>0\),
\[
 \sup_{\xi\in\Xi_n}Q_{n,\xi}^n
 \left(\left|\widehat V_n-V_{n,\xi}\right|>\epsilon\right)
 \longrightarrow0.                                                         \tag{36}
\]
Equations (33) and (36), together with \(V_{\mathrm{eff}}(P)>0\), prove uniform ratio consistency.  The distribution-function sandwich used pointwise, now applied along every sequence and combined with the same contradiction argument, proves uniform Wald coverage for \(\theta(Q_{n,\xi})\).

Finally, the assumed uniform likelihood expansion and mutual contiguity certify that \(Q_{n,\xi}^n\) are the declared contiguous local experiments and make base-law and local-law negligible events interchangeable.  The assumed uniform target expansion states
\[
 \sqrt n\{\theta(Q_{n,\xi})-\theta(P)\}
 =tE_P[\phi_{\mathrm{eff}}(P)s]+o(1)                                      \tag{37}
\]
uniformly over \(\Xi_n\).  Centering at \(\theta(P)\) therefore adds the shift in (37), whereas centering at the true local target removes it and leaves the common centered limit in (35).  This is regularity on every declared uniformly controlled bounded-score contiguous class.  At every local law \(\mathrm{thm:quotient-efficiency}\) makes \(\phi_{n,\xi}\) canonical, and (33) makes its variance converge to \(V_{\mathrm{eff}}(P)\).  Thus the estimator attains the local semiparametric efficiency bound, completing the proof.
